\scanpage{3}
\uline{Chap I}\qquad Introduction au Point de vue affine

\deleted{\uline{Point de vue affine}}

\textcircled{1} \uline{Définition \unsure{de|du} départ}

$E$ plan affine sur corps (disons \ldots) $k$. Un \uline{polygône}
\margin{\uline{Déf}}(déployé) dans $E$ est une partie finie \inserted{$S$} de $E$
(ensemble des sommets) \deleted{de cardinal $\geq 3$}, \uline{munie} d'une
``structure polygonale combinatoire'' i.e. d'une
partie à deux éléments $\omega=\{u,u^{-1}\}$ de $\unsure{\mathfrak{S}|\circlearrowright}_S$ qui soient
des permutations circulaires inverses l'une de l'autre (ce qui implique $n\geq3$). Les droites
joignant deux sommets contigus sont appelées les côtés\inserted{(en corr. 1-1 avec les côtés \unsure{comb})}\,; $\unsure{\omega|u}$ est \inserted{appelé} l'une des orientations de $P$. Si $n=\operatorname{card}S$,
c'est aussi le cardinal de l'ens des côtés -- on dit que $P$ est un $n$-polygône.

\margin{\unsure{pré-Déf|pre-Déf|re-Déf}}On dit que $P$ est \uline{régulier} si tout
autom combinatoire de \inserted{$Q=$} $(S,\omega)$ -- i.e. tout
élément du normalisateur \inserted{$G$} de $\omega$ dans $\mathfrak{S}_S$ --
se prolonge de façon unique en un autom
affine de $E$\,; il suffit bien sûr de l'exiger pour
\deleted{\unsure{les|des}} un système de gén, p.ex. \inserted{pr un $u\in\omega$} $u$ et une symétrie
$\sigma_s$ p.r. à un sommet $s$, ou une symétrie $\sigma_a$
p.r. à un \deleted{\unsure{côté}} \inserted{côté} $a$ [en fait, si $s$ et $a$ sont
incidents, deux \inserted{quelc} parmi les trois él. $\sigma_s,\sigma_a,u=\sigma_a\sigma_s$
forment \deleted{des} \inserted{un syst. de} générateurs]. Pour des raisons
techniques, on exige que les $u\in\omega$
(ou ce qui revient au $=$, les él. du
groupe \inserted{cyclique} $G^{\unsure{\bullet|0}}$ \unsure{engendré|engendrés} -- appelé \uline{groupe des
rotations} du polygône comb -- se réalisent
% archive sheet number: 2

